%==============================================================================
% Capitulo 9: Conclusiones
%==============================================================================
\section{Conclusiones}

El trabajo desarrollado ha cumplido satisfactoriamente con el objetivo general planteado en la Seccion~1.5: construir un sistema automatizado de mesa de ayuda capaz de recibir, procesar, clasificar y registrar incidentes de manera automática, reduciendo la intervención humana en la etapa inicial. La integración de N8N como motor de orquestación, un modulo Python con FastAPI como nucleo de procesamiento, PostgreSQL como capa de persistencia, Twilio como canal telefonico y el modelo Gemini~2.5 Flash como motor de inferencia linguistica ha demostrado ser una combinación arquitectonica viable, eficiente y económicamente accesible para organizaciones medianas.

\subsection{Cumplimiento de los objetivos del trabajo}

El diseñó arquitectonico del sistema materializo los principios de separación de responsabilidades y cohesion interna que la disciplina de ingenieria de software prescribe para sistemas distribuidos. La arquitectura de cinco capas ---canales de entrada, orquestación, procesamiento, inferencia y persistencia--- distribuyo las responsabilidades entre componentes independientes comunicados mediante interfaces REST, con TLS~1.3 a nivel del proxy inverso (Nginx) para la protección del perímetro del sistema. La eleccion de tecnologías maduras y de despliegue autoalojado (Docker, PostgreSQL, N8N) no fue meramente una decision instrumental sino una respuesta directa a las restricciones del contexto: organizaciones medianas argentinas que operan con presupuestos acotados, enfrentan limitaciones cambiarias para la adquisicion de software comercial y estan sujetas a un marco normativo que exige control sobre los datos personales.

Sobre esa arquitectura se construyó el clasificador hibrido, cuyo desempeño supero con holgura los umbrales establecidos en el segundo objetivo específico. La exactitud global del 92~\% y el F1 macro de 0,919 superan en siete y siete puntos porcentuales respectivamente las cotas minimas fijadas. Pero mas alla de la magnitud de las metricas, lo que los resultados revelan es la eficacia de la estrategia compositiva adoptada: combinar un filtro deterministico de primera etapa con un modelo de lenguaje grande invocado únicamente cuando el filtro no alcanza el umbral de confianza. Esta arquitectura de dos etapas produjo un sistema donde aproximadamente seis de cada diez incidentes se resuelven sin costo de inferencia externa, sin latencia de red y sin transferencia de datos a infraestructura de terceros. Los tres restantes, aquéllos que presentan ambiguedad semántica o redacción atipica, se benefician de la capacidad de comprension contextual del modelo de lenguaje. La distribucion del trabajo entre las dos etapas no fue un resultado buscado sino una propiedad emergente del diseñó, y sin embargo resultó ser uno de los hallazgos mas valiosos del trabajo: demuestra que los LLM, correctamente enmarcados como arbitros de último recurso y no como clasificadores universales, ofrecen un punto de equilibrio notable entre costo, latencia y exactitud.

La integración de los tres canales paralelos dentro de un único flujo de orquestación confirmó que la diversidad de vias de entrada no introduce dispersión operativa cuando la normalización se realiza tempranamente en el flujo. El canal telefonico, que a priori representaba el mayor desafio técnico por la inclusion de transcripción automática del habla, alcanzó una latencia de extremo a extremo de 12 a 15 segundos ---un valor que lo posiciona como una via prácticamente útil y no como una demostracion de laboratorio. La decision de normalizar la entrada en el propio N8N, antes de invocar al servicio de clasificación, resultó acertada: aislo la complejidad de cada protocolo de recepción del nucleo de procesamiento, permitiendo que el servicio Python operara siempre sobre la misma estructura de datos independientemente del canal de origen.

La evaluación comparativa entre el flujo manual y el automatizado arrojo resultados que trascienden lo meramente estadístico. La reducción de 165,3~s a 18,2~s en el tiempo medio de registro (\(p < 0{,}001\), \(r = 1{,}00\)) no deja espacio para la duda razonable: en las condiciones del experimento, el sistema automatizado fue invariablemente mas rapido que el operador humano, en los 200 pares medidos. La prueba de Wilcoxon con \(W = 0\) es, en si misma, un resultado poco frecuente en la investigación aplicada y merece ser interpretada con cautela: no implica que el sistema sea perfecto, sino que la diferencia con el proceso manual es tan abrumadora que ningun par invirtio el orden. La reducción de la intervención humana al 9,5~\% completa el cuadro: el sistema no solo es mas veloz sino que libera al personal técnico de la carga rutinaria de registro, redirigiendo ese tiempo hacía la resolución efectiva de los casos.

La documentación generada ---arquitectura, procedimientos de despliegue, especificación OpenAPI, flujo N8N exportable y análisis tipificado de errores--- constituye un activo que excede el cumplimiento del quinto objetivo específico. La inclusion de una tipologia de errores con tres patrones identificados (descripciones cortas, mezcla de categorías y jerga local) proporciona un punto de partida concreto para cualquier organización que desee replicar o adaptar el sistema, permitiendole anticipar las categorías de fallo mas probables antes de iniciar su propia implementación.

\subsection{Aportes y generalizacion}

Mas alla del cumplimiento de los objetivos específicos, este trabajo permite formular una afirmacion general de conocimiento relevante para el campo de la automatización de mesas de ayuda: los modelos de lenguaje grandes invocados mediante instrucciones contextualizadas ---sin ajuste fino sobre datos del dominio--- son capaces de alcanzar niveles de exactitud en clasificación de tickets de soporte superiores a los de clasificadores supervisados clasicos entrenados sobre representaciones estaticas, aun en idiomas de menor cobertura como el espanol rioplatense, siempre que se combinen con un filtro deterministico de primer nivel que absorba los casos triviales y reduzca el costo de inferencia.

Esta combinación arquitectonica ---orquestación visual, filtrado deterministico y LLM como arbitro semántico--- constituye un patron de diseñó reproducible para el dominio de gestión de servicios en organizaciones medianas de America Latina, donde la restricción presupuestaria y la soberanía de datos son determinantes para la adopción tecnologica. La arquitectura propuesta no depende de infraestructura especializada de GPU ni de licencias de software propietario, y su costo operativo esta dominado por el consumo de API de inferencia, que para el volumen observado (3.700 incidentes trimestrales) representa un costo marginal de aproximadamente 28 USD, una fraccion mínima de las 180 horas-persona liberadas por trimestre.

El trabajo aporta tres contribuciones que se sostienen mas alla del caso particular. La primera es una arquitectura reproducible documentada en todos sus niveles, desde el modelo de datos hasta el contrato REST y el flujo de orquestación, que cualquier equipo con conocimientos básicos de Docker, Python y N8N puede replicar. La segunda es evidencia empírica, recolectada bajo un diseñó experimental riguroso con doble etiquetado independiente y prueba no parametrica, sobre el desempeño de LLM en un dominio y una variante linguistica escasamente documentados en la literatura. La tercera es un marco de implementación que contempla desde el inicio las restricciones normativas argentinas, ofreciendo un camino de adopción que no requiere sortear barreras legales ni asumir riesgos de cumplimiento.

\newpage
